% Preface: the README's Om kursen, the course material and aids from info/README.md, and how the
% book is organised. No class, year or date: the book is meant to outlast them.

\chapter{Förord}
\markboth{Förord}{Förord}

Den här boken innehåller den matematik som används inom elektroteknik. Den är satt från en kurs i
elteknisk matematik, lektion för lektion, och behandlar:
\begin{itemize}
  \item Aritmetik, algebra samt linjära ekvationer och ekvationssystem.
  \item Potenser, rötter, andragradsekvationer samt exponentialfunktioner och logaritmer,
    inklusive decibel.
  \item Vektorer, funktioner och trigonometriska funktioner.
  \item Derivata och integraler, samt tillämpningar inom kretsteori.
  \item Komplexa tal på rektangulär, polär och eulerform, samt fasorer för
    växelströmsberäkningar.
\end{itemize}

\textbf{Matematiken används hela tiden på elektriska problem.} Redan i första kapitlet räknas
bråk på parallellkopplade motstånd, och boken återkommer ständigt till kretsarna:
\begin{itemize}
  \item Parallell- och seriekopplade motstånd, spänningsdelare samt effektberäkningar.
  \item Upp- och urladdning av kondensatorer i RC-kretsar.
  \item Växelspänningar: amplitud, frekvens, fas och sinusekvationer.
  \item Impedansberäkningar med komplexa tal och fasoraddition.
  \item Sampling av sinussignaler samt beräkning av derivator, integraler och komplexa tal i C,
    i laborationen i \kapref{18}.
\end{itemize}

\section*{Det här ska du kunna efter boken}
När du har arbetat dig igenom boken ska du kunna:
\begin{itemize}
  \item Förenkla algebraiska uttryck och lösa ekvationer och ekvationssystem som uppstår i
    eltekniska sammanhang.
  \item Räkna med potenser, rötter, logaritmer och decibel.
  \item Beräkna och tolka derivator och integraler för vanligt förekommande funktioner.
  \item Räkna med komplexa tal och fasorer för att analysera växelströmskretsar.
  \item Implementera kursens matematik i C med \file{math.h} och \file{complex.h}.
\end{itemize}

\section*{Vad du förväntas kunna}
Boken förutsätter grundskolans matematik. Det mesta av den repeteras ändå i de första kapitlen:
talmängder, räkneordning, bråk och procent i \kapref{1}, ekvationer i \kapref{3} och potenser i
\kapref{5}. Känner du dig osäker på grunderna är det alltså bara att börja från början.

Laborationen i \kapref{18} förutsätter dessutom grundläggande C-programmering: variabler,
funktioner, \code{printf()} och loopar. Flyttalen, \file{math.h} och \file{complex.h} gås igenom
i kapitlet.

\section*{Hur boken är organiserad}
Varje kapitel motsvarar en lektion i kursen, med samma nummer: kapitel~5 är lektion~5, och dess
material finns i kursens repository under \file{lectures/L05}.

\begin{booktable}{@{}p{5em}L@{}}
  \thead{Kapitel} & \thead{Ämne} \\
  \midrule
  1 & Talsystem och grundläggande aritmetik \\
  2 & Algebra \\
  3 & Linjära ekvationer och olikheter \\
  4 & Ekvationssystem \\
  5 & Potenser och rötter \\
  6 & Andragradsekvationer, med övningsdugga 1 \\
  7 & Vektorer \\
  8 & Funktioner \\
  9 & Trigonometriska funktioner \\
  10 & Exponentialfunktioner, logaritmer och decibel \\
  11 & Repetition av kapitel 7--10, med övningsdugga 2 \\
  12 & Derivata, del I \\
  13 & Derivata, del II \\
  14 & Integraler \\
  15 & Komplexa tal, del I \\
  16 & Komplexa tal, del II \\
  17 & Komplexa tal, del III \\
  18 & Matematik i C: laboration \\
  19 & Repetition och tentamensförberedelse, med övningstentamen \\
\end{booktable}

Ett kapitel börjar med vad det innehåller och går sedan igenom teorin, med exempel som visar
varje steg. Därefter kommer en sammanfattning med kapitlets formler, vad du ska kunna och några
frågor under rubriken \emph{Testa dig själv}, och sist kapitlets uppgifter. Kapitel~11 och~19 har
ingen ny teori; de repeterar och innehåller i stället en övningsdugga respektive en
övningstentamen.

\section*{Uppgifterna}
Uppgifterna i varje kapitel är indelade i delar och numrerade som i kursen, så att uppgift~2.3 är
den tredje uppgiften i del~2:

\begin{booktable}{@{}p{5em}L@{}}
  \thead{Del} & \thead{Vad den innehåller} \\
  \midrule
  Del 1 & Inledande uppgifter, eller repetition av föregående kapitel. \\
  Del 2 & Nytt stoff: kapitlets egen teori, ofta tillämpad på en krets. \\
  Del 3 & Extrauppgifter att träna vidare på efter lektionen. \\
\end{booktable}

I kursen räknas del~1 och del~2 under lektionen, först på egen hand och sedan gemensamt, medan
del~3 är frivillig träning.

\section*{Facit och lösningsförslag}
\textbf{Svaren på samtliga uppgifter finns i facit längst bak i boken.} Varje uppgift har en liten
etikett, \svarexempel, som anger på vilken sida dess svar står; i PDF-versionen är både
etiketten och uppgiftens nummer i facit länkar mellan uppgift och svar.

\textbf{Facit innehåller bara svaren, inte lösningarna.} Ett svar visar om du har räknat rätt,
men inte hur man kommer fram till det. Fullständiga lösningsförslag publiceras i kursens
repository efter respektive lektion, under \file{lectures/LNN/appendix/c_solutions.md}. Räkna
alltid själv först, jämför sedan med facit, och läs lösningsförslaget när svaret inte stämmer
eller när du vill se en annan väg fram.

\section*{Hjälpmedel}
En miniräknare, vilken som helst, används genom hela boken. Många uppgifter är ändå tänkta att
räknas utan den, och de säger det. Laborationen i \kapref{18} kräver bara en webbläsare: koden
skrivs och körs i en kompilator på webben, utan installation.

\section*{Beteckningar}
Boken följer svensk skrivpraxis och de beteckningar som är vanliga inom elektroteknik:
\begin{itemize}
  \item \textbf{Decimalkomma}: $4{,}7$, inte $4.7$. Undantaget är C-koden i \kapref{18}, där
    decimaltecknet är en punkt.
  \item \textbf{Enheter} skrivs med ett litet mellanrum efter talet: $4{,}7\,\text{k}\Omega$,
    $50\,\text{Hz}$, $30\,\%$.
  \item \textbf{Resistans skrivs före ström}, som i Ohms lag: $U = RI$ och $P = RI^2$.
  \item \textbf{Parallellkoppling} skrivs $R_1 \parallell R_2$.
  \item \textbf{Logaritmer}: $\log$ är tiologaritmen och $\ln$ den naturliga logaritmen.
  \item \textbf{Den imaginära enheten} skrivs $j$, eftersom $i$ redan betecknar ström.
  \item \textbf{Vinklar} anges i radianer om inget annat sägs; grader markeras med $^{\circ}$.
\end{itemize}
